\subsection{Antifascist Ethics for AI Alignment}
\label{sec:2.1.4}
Antifascist ethics opposes concentrating power and authority in AI systems and counters the potential misuse of AI technology by authoritarian actors. The word is chosen deliberately: fascism is the narrower and more charged term, and I use it because it is more honest about the politics in which this book is being written — but the target is authoritarian power in every form, and a system built on these principles is anti-authoritarian in full, not only antifascist.

A stance needs a definition to be a design specification, and this is where books of this kind usually stop. Opposing “concentrated, unaccountable power” gives an engineer nothing to build. It does not say what a system would detect, what would count as a false positive, or how anyone would know the system was wrong. So before anything a designer could be asked to build, the definition.

\runin{Fascism as a structural signature}
I take fascism to name a structure rather than a costume. Four features travel together, and each is visible from outside the institution displaying it:

\begin{enumerate}
  \item Recuperation of dissent. The institution does not suppress disagreement; it metabolizes it. Refusal is converted into a managed appearance — a permitted outlet, a consultation, a feedback channel — which the institution can then point to as proof that it listens. Debord aimed this warning at his own concept, noting that even denunciation of the spectacle can become a hollow formula that reinforces what it denounces, and that dissatisfaction itself can be sold back as a commodity of opposition \autocite{debord1967spectacle}. The tell is where what the mechanism carries stops. The channel works, and the institution uses what comes through it; what has nowhere to go is the judgment that the frame the channel operates inside is the wrong frame.
  \item Aestheticization of the metric. The measure substitutes for the thing measured, and the substitution is experienced as rigor. Benjamin's account of fascism organizing the masses by granting expression in place of rights is the political form of this \autocite{benjamin1935artwork}; the administrative form is a number that collapses two states — done and verified — into one, so that whatever falls outside the number is not refuted or dismissed but simply unnamed.
  \item Coding of exception as betrayal. Marking a case as one the rule does not fit becomes an act requiring justification, and then an act of disloyalty. The honest statement that a standard does not apply to this case reads as weakness. It is in fact the irreducible labor of deciding, case by case, what a rule is for.
  \item Decoupling of the model from the world it claims to track. The institution's representation of reality and reality itself come apart, and the institution defends the representation. This is the load-bearing one, because the other three are what protect the decoupling from correction.
\end{enumerate}
Those four are testable. A system asked whether an institution shows them is being asked about observable artifacts: whether the dissent channel has ever changed an outcome, whether the metric admits an unmeasured remainder, what happens to the person who files the exception, whether the model is ever checked against anything outside itself.

The first of the four is the one this book meets again in its own technical material, twice, at sizes that look unrelated: in how a training corpus gets labeled, and in how a model is tuned on human feedback. Both arrangements collect disagreement, put it to work, and are built so that none of it can reach the scheme it disagrees with. Chapter~\ref{sec:7} puts the two instances together and works out what follows, including for the argument being made here.

\runin{Two scales}
Deleuze and Guattari's distinction between molar and molecular fascism is the reason the four features recur at sizes that seem unrelated \autocite{deleuze1987plateaus}. Molar fascism is the state, the regime, the camp. Molecular fascism is the office someone commutes to, the form they fill out, and the desire with which they fill it out. The same administrative pattern is being rediscovered at each scale, which is why a definition pitched only at the state cannot see the version an AI deployment actually instantiates. At state scale the pattern kills journalists. At organizational scale it turns an engineering decision into scripture: the constraint remains, its reason dies, and asking why becomes the deviant act. A system that can only recognize the molar case will certify the molecular one as healthy.

\runin{What makes it fascism and not authoritarianism in general}
The structural signature is shared with a good deal of ordinary institutional decay, and a detector tuned to it alone will flag every badly run organization on earth. Fascism proper adds features that are narrower and, for the same reason, more discriminating: mass mobilization, which mere obedience does not require; a leader cult; a scapegoated out-group, which Eco places at the origin — the first appeal is against the intruder \autocite{eco1995urfascism}; contempt for procedural legality as such, so that law is an obstacle to the national will and not a constraint on power; and deniability as technique, in which the movement's violence is disowned at the top while being rewarded at the bottom. Paxton's compressed version is worth having in full: fascism is “a form of political behavior marked by obsessive preoccupation with community decline, humiliation, or victimhood and by compensatory cults of unity, energy, and purity” \autocite{paxton2004anatomy}.

Those are detection targets. Concentrated, unaccountable power is not.

\runin{The scale problem, and what this book does about it}
Put the two halves of this section together and they do not fit. Every discriminating feature is a property of a mass political movement: mobilization needs a crowd, a leader cult needs a leader, a scapegoat needs a public. An AI deployment has a form, a threshold, a review queue, an escalation path. It operates at the molecular scale, and that is the only scale it has. There the discriminators as stated are unavailable, which appears to leave the four-feature signature alone — and the four-feature signature, as I have just said, flags every badly run organization on earth.

I take that as a finding, and it is the more demanding reading. Molecular fascism is the thing itself at the size where it is actually lived, and not an early warning of a regime that has yet to assemble. The office that converts refusal into a feedback form does the same work whether or not a regime ever assembles out of a thousand offices like it. A detector firing on the four features across a thousand institutions is reporting its target.

Saying that is not yet an argument, and introducing a distinction does not license it. Deleuze and Guattari supply the vocabulary and no authority beyond it, so the work has to be done here. Two things do it.

The first is that the molar case has no mechanism of its own. A regime does nothing except through offices, forms, thresholds, escalation paths, and people who fill them out with some degree of enthusiasm. When a state kills journalists, the killing is authorized on a form, justified in a category, and carried out by someone whose promotion depends on it. What assembling adds is scale, protection from consequence, and the removal of the outside bodies that might have interrupted any single office. It does not add a new kind of thing. If the mechanism at the large scale is the small-scale mechanism repeated and shielded, then the name belongs to the mechanism. The natural objection is that extent is not nothing — a regime that disappears people and an office that loses paperwork are not equivalent, and no argument should imply they are. Agreed, and the claim is about kind rather than severity. A definition is not a sentencing guideline.

The second does more work. The discriminators are unavailable at the molecular scale only in the form the previous paragraphs stated them, and each has a molecular restatement that is not a property of ordinary decay. Mass mobilization becomes participation exceeding what the rule requires: staff who volunteer for the disfavored task, elaborate on it, take pride in it. A leader cult becomes personal loyalty overriding role — exceptions granted by proximity to someone rather than by procedure, and a decision reversed because of who asked. The scapegoat becomes an internal category defined so that adverse treatment of it is the category's purpose: the flagged applicant, the difficult claimant, the file everyone knows is handled differently. Contempt for procedural legality becomes procedure enforced downward and waived upward. And deniability as technique keeps its shape exactly: the practice is disowned in the policy and rewarded in the promotion.

That is what separates fascism from decay at the size an AI deployment operates at. Decay is an institution failing to do what it says. This is an institution succeeding at something else, with participation, a target, and a story about why the rules do not apply to whoever is applying them. The four features tell you a place is sick; these five tell you what it is sick with. I am not claiming they are easy to observe — they are harder than the four, since each requires knowing what was rewarded rather than what was written — and section~\ref{sec:11.3} is where that difficulty gets its due.

This widens what the word covers. On this reading “antifascist” names the political inheritance of a structural claim: where the analysis comes from, and what it was built to fight. It does not mark off a narrow class of institutions that a classifier separates from ordinary decay. The molar discriminators keep their own job, which is to tell a historian, a prosecutor, or a citizen that a movement has crossed into territory ordinary politics does not contain. No system is being asked to make that call, and this book does not ask one to. A reader who thinks the word belongs to the molar case will think I have diluted it, and the answer is the two arguments above rather than the distinction itself: the mechanism is one mechanism, and the discriminating features have molecular forms that ordinary institutional failure does not exhibit.

That a detector of these four features would find something nearly everywhere is the expected result and not a defect of the instrument. The tendency is permanent; what varies is how far it has got in a particular place, and whether it is getting further. Which is why the question worth asking of an institution is not whether the signature is present but how much of it is, and in which direction it has moved since anyone last looked — section~\ref{sec:9.3.4}'s slope rather than its level, applied to the definition itself and not only to a dissent channel.

One consequence travels with the reader through everything that follows, because the design chapters assume a capability this framing has widened. A tool reporting the four features will have something to report nearly everywhere, and its output is an accusation. In the hands of whoever holds it, that is a technical-sounding case that any given opponent is operating fascistically: a weapon handed to the actor it was built to constrain, and the dual-use problem of section~\ref{sec:5.4.3} at its sharpest, since what comes out of it is a charge. The opposite failure is no safer. A tool tuned until it reports nothing supplies false assurance, and an absence of warnings reads as evidence of health. Better engineering fixes neither. Both are governed by what the tool may output and to whom, which makes them a constraint the specification has to carry rather than a setting chosen at deployment.

What stays genuinely unsolved is the molecular case itself: telling a recuperated dissent mechanism from a merely mediocre one, when both produce the same paperwork. Section~\ref{sec:11.3} takes that up, along with what a detector would need that this book cannot yet tell an engineer how to obtain. Naming the targets was the easy half.

In AI alignment terms, then, antifascist design means building systems that enhance democratic values, civil rights, and inclusive solidarity, that refuse to facilitate autocratic practices or prop up an oppressive regime, and that do not themselves display the four features above. What that comes to in practice is the subject of the rest of the book and not of this section: decision procedures that answer to democratic values, privacy and autonomy built in by construction and not left to policy, power and resources distributed instead of concentrated, an explanation a person can actually read and an independent body positioned to audit it, architectures that resist the concentration of the field in a handful of firms, misuse caught while it is happening, and systems built to widen a conversation instead of manufacturing a feeling of agreement inside it. Section~\ref{sec:6.4.2} and chapters~\ref{sec:8} and \ref{sec:9} are where each of those is worked out and priced.
